\clearpage
\setcounter{aufgabe}{15}
\einheitskopf{Einheit 3 von 5 · Prozentwert berechnen}\label{e3}

\begin{aufgabe}{Prozentwert – Streifen einteilen. Der ganze Streifen sind 100\,\%. Teile ihn mit dem Lineal ein und färbe den verlangten Teil grau. Rechne noch nicht.}
\begin{teile}
\teil Teile in 10-\%-Schritte ein und färbe 30\,\% grau.
\teil Teile in 25-\%-Schritte ein und färbe 75\,\% grau.
\teil Teile in 10-\%-Schritte ein und färbe 5\,\% grau.
\teil Färbe ohne Einteilung ungefähr 60\,\% grau.
\end{teile}
\streifenreihe[0]{a/0,b/0,c/0,d/0}
\end{aufgabe}

\begin{aufgabe}{Prozentwert – am Streifen. Der ganze Streifen ist das Ganze. Wie groß ist der graue Teil? Schreibe die Einheit dazu. Beispiel: Sind 2 der 10 Kästchen eines Streifens von $50\,€$ grau, so sind das $10\,€$.}
\streifenwertreihe{a/50/?/80 €, b/25/?/60 kg, c/10/?/200 m, d/75/?/24 l}
\begin{teilezwei}
\tz \leerfeld & \tz \leerfeld \\
\tz \leerfeld & \tz \leerfeld \\
\end{teilezwei}
\end{aufgabe}

\begin{aufgabe}{Prozentwert – rechnen. Rechne mit dem Ein-Prozent-Weg: erst das Ganze durch 100, dann mal dem Prozentsatz.}
\beispiel{\text{8\,\% von 300\,kg:}\quad 1\,\% &= 300\,\text{kg} : 100 = 3\,\text{kg} \\ 8\,\% &= 8 \cdot 3\,\text{kg} = 24\,\text{kg}}
\begin{teilezwei}
\tz $50\,\%$ von $60\,€$: \leerfeld & \tz $25\,\%$ von $80\,$kg: \leerfeld \\
\tz $10\,\%$ von $70\,$m: \leerfeld & \tz $50\,\%$ von $46\,€$: \leerfeld \\
\tz $7\,\%$ von $300\,€$: \leerfeld & \tz $15\,\%$ von $80\,€$: \leerfeld \\
\end{teilezwei}
Rechne hier kurz: schreibe den Prozentsatz als Dezimalzahl und nimm damit das Ganze mal.
\begin{teile}
\teil $35\,\%$ von $140\,$kg: \leerfeld
\teil $40\,\%$ von $7{,}50\,€$: \leerfeld
\teil $120\,\%$ von $50\,$kg: \leerfeld
\end{teile}
\end{aufgabe}

\begin{aufgabe}{Prozentwert – aus dem Sachtext. Lies die Frage genau: Manchmal ist die Ersparnis gesucht, manchmal der neue Preis.}
\begin{teile}
\teil Ein Rucksack kostet $45\,€$. Im Schlussverkauf gibt es $20\,\%$ Rabatt. Wie viel Euro spart man? \leerfeld[€]
\teil Wie viel muss man für den Rucksack dann bezahlen? \leerfeld[€]
\teil Ein Handwerker berechnet $240\,€$ für seine Arbeit. Dazu kommen $19\,\%$ Mehrwertsteuer. Wie viel Euro Mehrwertsteuer sind das? \leerfeld[€]
\teil Ein Zelt kostet $320\,€$ und wird um $15\,\%$ billiger. Wie viel Euro spart man? Kreuze an. \hfill (P10 2021) \\
\kreuz{$272\,€$}\\
\kreuz{$48\,€$}\\
\kreuz{$21{,}33\,€$}
\end{teile}
\end{aufgabe}

\begin{aufgabe}{Prozentwert – Fehler finden. Lena soll $40\,\%$ von $60\,€$ berechnen. Sie rechnet:}
\rechnung{0{,}04 \cdot 60\,€ &= 2{,}40\,€}
\begin{teile}
\teil Benenne den Fehler und korrigiere die Rechnung. \\ Fehler: \dsleer\dsleer\dsleer \quad richtig: \leerfeld[€]
\teil Berechne $70\,\%$ von $90\,€$. \leerfeld[€]
\end{teile}
\end{aufgabe}

\begin{aufgabe}{Prozentwert – Begründen. Begründe in einem Satz.}
\begin{teile}
\teil Warum teilt man beim Ein-Prozent-Weg zuerst das Ganze durch 100? \\ \dsleer\dsleer\dsleer\dsleer
\teil Jan sagt: „$25\,\%$ von $80\,€$ ist genauso viel wie $80\,\%$ von $25\,€$.“ Hat er recht? \janein \\ Begründung: \dsleer\dsleer\dsleer\dsleer
\end{teile}
\end{aufgabe}

\begin{aufgabe}{Prozentwert – Anwendung. In Altstadt leben 24\,000 Menschen, davon sind $45\,\%$ jünger als 30 Jahre. In Neustadt leben 15\,000 Menschen, davon sind $60\,\%$ jünger als 30 Jahre.}
\begin{teile}
\teil Wie viele Menschen unter 30 Jahren leben in Altstadt? \leerfeld
\teil Wie viele Menschen unter 30 Jahren leben in Neustadt? \leerfeld
\teil In welcher Stadt leben mehr Menschen unter 30 Jahren? Erkläre, warum der größere Prozentsatz nicht die größere Anzahl bedeutet. \\ \dsleer\dsleer\dsleer\dsleer\dsleer
\end{teile}
\end{aufgabe}
